\section{Introduction}\label{sec:intro}

Person detection is the entry point of many video-analytics applications: occupancy counting, safety monitoring, footfall analysis and smart-building control. One-stage detectors of the YOLO family~\cite{redmon2016yolo,jocher2023yolov8} deliver accurate boxes in a single forward pass, but that pass is expensive on the CPUs and embedded boards on which camera streams are often processed. On one CPU thread of our test machine, YOLOv8m needs \YoloMs\,ms per $640\times640$ frame, more than ten times the frame period of a 30\,FPS stream. At the same time, a static camera observes the same background for most of the day, so running the detector on every frame mostly recomputes an answer that has not changed.

The obvious remedy is a \emph{motion gate}: a cheap change detector inspects every frame, the expensive detector runs only when change is reported, and the last detections are reused otherwise. Frame differencing~\cite{jain1979difference} or background subtraction~\cite{stauffer1999adaptive,zivkovic2004improved} costs about a millisecond per frame, and difference-based frame filtering is a core ingredient of video-analytics systems~\cite{kang2017noscope,li2020reducto,arefeen2022framehopper}. Gating is not free, however. It introduces a new kind of error, \emph{staleness}: reused boxes lag behind the scene. And the gate can fail in two directions. False triggers caused by sensor noise, camera jitter or illumination changes waste detector calls and erode the savings; missed triggers caused by small or slow objects, or by a person leaving the image, leave outdated boxes.

This paper revisits a prototype of such a hybrid system, which gated YOLOv8m with OpenCV frame differencing and reported a $1.71\times$ speed-up on a small synthetic benchmark. In that benchmark the gated and the ungated system were both perfect, the sensitivity ablation of the gate was flat, and the one scenario in which the gate failed (camera jitter) remained unexplained, so the evaluation could not say \emph{when} gating is safe. We turn the prototype into an analysed and tested design and make four contributions:

\begin{itemize}
\item \textbf{A scheduler with guarantees.} Two interpretable parameters, a refresh interval $K$ and a minimum interval $M$, bound the age of every reported detection by $K-1$ frames and the fraction of frames sent to the detector between $1/K$ and $1/M$, whatever the scene (\cref{sec:scheduler}). A cost model yields the speed-up and the break-even condition and is validated end to end.
\item \textbf{An analytical failure envelope of frame differencing.} The original gate takes absolute differences \emph{before} smoothing. Noise therefore contributes a mean of $\BiasCoef\,\sigma$ that smoothing cannot remove, a \emph{rectification bias}; combined with a percolation argument, it predicts the noise level at which the gate breaks down. Simple models predict the response to illumination steps and to small or slow objects (\cref{sec:envelope}). All predictions are tested experimentally.
\item \textbf{A stabilised gate (\FDS).} Smoothing a \emph{signed} difference, registering consecutive frames by phase correlation guarded by an edge-based model selection, compensating global intensity offsets and propagating the registration error into the threshold removes most false triggers caused by noise, illumination changes and camera jitter, at a cost of a few milliseconds per frame (\cref{sec:fds}).
\item \textbf{An evaluation protocol and an open implementation.} A controlled benchmark with exact ground truth isolates nuisance factors one at a time, every policy is compared with fixed-rate subsampling at the same computational budget, and a real video serves as a held-out check. The Python package, a command-line tool, unit tests and a single script that regenerates every number, table and figure of this paper accompany it.
\end{itemize}

\noindent From a practical point of view, the analysis leads to simple deployment rules (\cref{sec:discussion}): choose $K$ from the staleness the application tolerates, use $M$ to cap the cost in busy scenes, and use the stabilised gate whenever the camera may vibrate or the lighting may change.
